\section{Preliminaries}\label{sec:preliminaries}

\subsection{Symmetric functions and \texorpdfstring{$P$}{P}-partitions}

We use symmetric functions in $\xvec=(x_1,x_2,\dotsc)$ with coefficients in
$\setZ[q]$, see~\cite{Stanley1972} for the background on $P$-partitions.
For an integer partition $\lambda$ we write
$\elementaryE_\lambda=\elementaryE_{\lambda_1}\elementaryE_{\lambda_2}\dotsm$.
For a map $\kappa$ from a finite set $V$ to $\setNN$ we put
$x_\kappa\coloneqq\prod_{u\in V}x_{\kappa(u)}$.

Let $P$ be a partial order on a finite set $V$.
A map $\kappa:V\to\setNN$ is a \defin{strict $P$-partition} if
$\kappa(u)<\kappa(v)$ whenever $u<_P v$, and we set
\[
 \defin{\Kp_P(\xvec)}\coloneqq\sum_{\kappa}x_\kappa,
\]
the sum ranging over all strict $P$-partitions.
This is a quasisymmetric function, and it is symmetric when $P$ is a disjoint
union of chains: a chain with $k$ elements contributes $\elementaryE_k$.
We write $u\lessdot_P v$ when $v$ \defin{covers} $u$ in $P$.
If $D$ is a directed graph whose transitive closure is a partial order, then
$\Kp_P$ for this order counts the maps that increase strictly along every edge
of $D$, and every cover of the order is an edge of $D$.

\subsection{Decorated unit-interval graphs}

An \defin{area sequence} is a tuple $\avec=(a_1,\dotsc,a_n)$ of nonnegative
integers with $a_1=0$ and $a_{i+1}\leq a_i+1$.
We put $\defin{b(i)}\coloneqq i-a_i$, so that $b$ is weakly increasing.
A pair $uv$ with $u<v$ is a \defin{natural edge} if $b(v)\leq u$.
The natural edges form a unit-interval graph: if $uw$ is a natural edge and
$u<v<w$, then $uv$ and $vw$ are natural edges.
There are $a_v$ natural edges ending at $v$, so $\sum_i a_i$ natural edges in
total.

A natural edge $uv$ is a \defin{corner} if no other natural edge $u'v'$
satisfies $u'\leq u$ and $v\leq v'$.
Equivalently, $u=b(v)$, and either $v=n$ or $b(v+1)>b(v)$.
Every natural edge lies inside a corner, and distinct corners have distinct
tails and distinct heads.
A \defin{decorated unit-interval graph} $\Gamma=(\avec,S)$ is an area sequence
together with a set $S$ of corners, the \defin{strict edges}.
The other natural edges are \defin{normal}.
Decorated unit-interval graphs are the graphs $\Gamma_P$ of Schröder paths $P$
in~\cite[Sec.~2.6]{AlexanderssonSulzgruber2022}: the natural edges are the
cells below the Dyck path obtained by replacing every diagonal step by a north
step followed by an east step, and the strict edges are the cells of the
diagonal steps.

We draw $\Gamma$ as in~\cite{AlexanderssonSulzgruber2022}.
The cell in column $u$ and row $v$ represents the pair $uv$; white cells are
normal edges, shaded cells are not natural edges, and a shaded cell with an
arrow is a strict edge.

\begin{example}\label{ex:running}
The Schröder path $\nstep\nstep\nstep\dstep\dstep\estep\nstep\estep\estep\estep$
of~\cite[Ex.~2.8]{AlexanderssonSulzgruber2022} gives the area sequence
$(0,1,2,3,3,2)$ with strict edges $14$ and $25$.
The corners are $14$, $25$ and $46$; only the first two are strict.
In the left diagram below, arrows in white cells mark the ascending edges of an
orientation $\theta$, defined in the next subsection:
\[
\ytableausetup{boxsize=1.1em}
\begin{ytableau}
*(lightgray) & *(lightgray) & *(lightgray) & & \to & *(yellow) 6 \\
*(lightgray) & *(lightgray) \to & & & *(yellow) 5 \\
*(lightgray) \to & \to & \to & *(yellow) 4 \\
\to & \to & *(yellow) 3 \\
 & *(yellow) 2 \\
*(yellow) 1
\end{ytableau}
\qquad\qquad
\begin{tikzpicture}[x=2.4em,y=2.4em]
\node[vtx] (1) at (0,0) {1};
\node[vtx] (3) at (0,1) {3};
\node[vtx] (4) at (0,2) {4};
\node[vtx] (2) at (1.5,0.5) {2};
\node[vtx] (5) at (1.5,1.5) {5};
\node[vtx] (6) at (1.5,2.5) {6};
\draw[sel] (1)--(3);
\draw[sel] (3)--(4);
\draw[strictedge] (1) to[bend left=40] (4);
\draw[sel] (2)--(3);
\draw[sel] (2)--(4);
\draw[strictedge] (2)--(5);
\draw[sel] (5)--(6);
\end{tikzpicture}
\]
The right picture shows the strict edges (thick) and ascending edges of
$\theta$.
\end{example}

\subsection{Colorings and orientations}

A \defin{coloring} of $\Gamma$ is a map $\kappa:[n]\to\setNN$ with
$\kappa(u)<\kappa(v)$ for every strict edge $uv$.
An \defin{ascent} of $\kappa$ is a normal edge $uv$ with $\kappa(u)<\kappa(v)$,
and $\asc(\kappa)$ is the number of ascents.
The \defin{vertical-strip LLT polynomial} is
\[
 \defin{\LLT_\Gamma(\xvec;q)}\coloneqq\sum_{\kappa}q^{\asc(\kappa)}x_\kappa,
\]
the sum ranging over all colorings of $\Gamma$.
It is a symmetric function~\cite{AlexanderssonPanova2018,HaglundHaimanLoehr2005},
and every vertical-strip LLT polynomial of~\cite{LLT1997} arises in this
way~\cite{AlexanderssonPanova2018,CarlssonMellit2017}.

An \defin{orientation} $\theta$ of $\Gamma$ directs every normal edge.
A normal edge $uv$ with $u<v$ is \defin{ascending} in $\theta$ if it is directed
from $u$ to $v$, and $\asc(\theta)$ is the number of ascending edges.
Strict edges are always directed upwards.
We write $\ORIENTS(\Gamma)$ for the set of orientations, and we identify
$\theta$ with its set of ascending edges.
Let $\Dplus{\theta}$ be the directed graph on $[n]$ whose edges are the strict
edges and the ascending edges of $\theta$.
All its edges point upwards, so its transitive closure is a partial order
$\Pplus{\theta}$ on $[n]$.

\begin{lemma}\label{lem:stateExpansion}
For every decorated unit-interval graph $\Gamma$,
\[
 \LLT_\Gamma(\xvec;q+1)=\sum_{\theta\in\ORIENTS(\Gamma)}q^{\asc(\theta)}\,
 \Kp_{\Pplus{\theta}}(\xvec).
\]
\end{lemma}

\begin{proof}
For a coloring $\kappa$, expand $(q+1)^{\asc(\kappa)}$ as a sum of $q^{|A|}$
over the subsets $A$ of the ascents of $\kappa$.
A pair $(\kappa,A)$ is the same as an orientation $\theta=A$ together with a
strict $\Pplus{\theta}$-partition $\kappa$, since $\kappa$ increases on the
strict edges and on the edges of $A$.
The claim follows.
\end{proof}

The \defin{highest reachable vertex} $\hrv{\theta}{u}$ of a vertex $u$ is the
largest $v$ with $u\leq v$ in $\Pplus{\theta}$.
The fibres of $\hrv{\theta}{\cdot}$ form a set partition of $[n]$, and
$\defin{\lambda(\theta)}$ is the partition given by their sizes.

\begin{example}
For the orientation $\theta$ in Example~\ref{ex:running}, the vertices
$2,5,6$ reach $6$, and the vertices $1,3,4$ reach $4$.
Thus $\lambda(\theta)=(3,3)$ and $\asc(\theta)=5$, see
also~\cite[Ex.~2.8]{AlexanderssonSulzgruber2022}.
\end{example}

The following expansion was conjectured
in~\cite{GarsiaHaglundQiuRomero2019,Alexandersson2019llt}.
Its positivity was obtained independently by M.~D'Adderio~\cite{DAdderio2020}.

\begin{theorem}[{\cite[Cor.~2.10]{AlexanderssonSulzgruber2022}}]
\label{thm:AS}
For every decorated unit-interval graph $\Gamma$,
\[
 \LLT_\Gamma(\xvec;q+1)=\sum_{\theta\in\ORIENTS(\Gamma)}q^{\asc(\theta)}\,
 \elementaryE_{\lambda(\theta)}(\xvec).
\]
\end{theorem}

Comparing with Lemma~\ref{lem:stateExpansion}, the function
$\Kp_{\Pplus{\theta}}$ and the elementary function $\elementaryE_{\lambda(\theta)}$
agree after summing over all orientations with a given number of ascending
edges, although they differ for individual orientations.
The purpose of this paper is to identify much smaller sets of orientations over
which they agree.

\subsection{Flips}

\begin{definition}\label{def:flip}
Let $u<v<w$ be vertices such that $uw$ is a natural edge, and let $\theta$ be an
orientation in which $uv$ is ascending and $vw$ is not.
Let $\theta'$ be obtained from $\theta$ by reversing both $uv$ and $vw$, so that
$vw$ is ascending in $\theta'$ and $uv$ is not.
We say that $\theta$ and $\theta'$ differ by a \defin{flip} if
\[
 u\lessdot v \text{ and } u\lessdot w \text{ in } \Pplus{\theta},
 \qquad\text{and}\qquad
 v\lessdot w \text{ and } u\lessdot w \text{ in } \Pplus{\theta'}.
\]
The \defin{flip classes} of $\Gamma$ are the equivalence classes of
$\ORIENTS(\Gamma)$ under the relation generated by flips.
\end{definition}

Since $u\lessdot w$ in $\Pplus{\theta}$, the pair $uw$ is an edge of
$\Dplus{\theta}$, the \defin{long edge} of the flip.
It is retained, and so are all other edges of $\Dplus{\theta}$.
The two exchanged \defin{short edges} $uv$ and $vw$ lie strictly inside $uw$;
hence they are not corners, and in particular they are normal.
The definition is symmetric in $\theta$ and $\theta'$, with the roles of the two
short edges exchanged.
It is the LLT counterpart of the flips introduced by B.-H.~Hwang for
acyclic orientations~\cite{Hwang2022x}, which reverse a fork of two
covers.

In an ordinary graph, two of the four cover conditions are automatic.

\begin{lemma}\label{lem:flipCriterion}
Let $\Gamma$ be ordinary, let $u<v<w$ with $uw$ natural, and let $\theta$ be an
orientation in which $uv$ and $uw$ are covers of $\Pplus{\theta}$.
Then $\theta$ and $\theta-uv+vw$ differ by a flip.
Similarly, if $vw$ and $uw$ are covers of $\Pplus{\theta}$, then $\theta$ and
$\theta-vw+uv$ differ by a flip.
\end{lemma}

\begin{proof}
For the first claim, $vw\notin\theta$ since $uw$ is a cover; put
$\theta'=\theta-uv+vw$.
A path from $v$ to $w$ in $\Dplus{\theta'}$ other than the edge $vw$ does not use
$uv$, so it is a path in $\Dplus{\theta}$, and preceded by $uv$ it would bypass
$uw$.
A path of length at least two from $u$ to $w$ in $\Dplus{\theta'}$ either avoids
$vw$, and then bypasses $uw$ in $\Dplus{\theta}$, or consists of a path from $u$
to $v$ followed by $vw$; that path from $u$ to $v$ avoids $uv$ and $vw$, so it
bypasses $uv$ in $\Dplus{\theta}$.
Hence $v\lessdot w$ and $u\lessdot w$ in $\Pplus{\theta'}$.
The second claim follows in the same way.
\end{proof}

\begin{example}\label{ex:smallestFlip}
Let $\Gamma$ be the complete graph on $[3]$, with area sequence $(0,1,2)$ and no
strict edges.
The two orientations
\[
\ytableausetup{boxsize=1.1em}
\theta=\;
\begin{ytableau}
\to & & *(yellow) 3 \\
\to & *(yellow) 2 \\
*(yellow) 1
\end{ytableau}
\qquad\text{and}\qquad
\theta'=\;
\begin{ytableau}
\to & \to & *(yellow) 3 \\
 & *(yellow) 2 \\
*(yellow) 1
\end{ytableau}
\]
differ by a flip at $(1,2,3)$:
the covers of $\Pplus{\theta}$ are $1\lessdot 2$ and $1\lessdot 3$, and those of
$\Pplus{\theta'}$ are $1\lessdot3$ and $2\lessdot3$.
No other flip applies to them, so $C=\{\theta,\theta'\}$ is a flip class.
The vertex $2$ is its own highest reachable vertex in $\theta$, while all
vertices reach $3$ in $\theta'$.
Hence
\[
 \Kp_{\Pplus{\theta}}+\Kp_{\Pplus{\theta'}}
 =\elementaryE_{21}+\elementaryE_3,
\]
which can be checked on the coefficients of $x_1x_2x_3$ and $x_1^2x_2$.
Neither summand on the left is symmetric.
\end{example}

\begin{lemma}\label{lem:flipInvariants}
On a flip class, the number of ascending edges, the connected components of the
underlying undirected graph of $\Dplus{\theta}$, and the set of corners which
are edges of $\Dplus{\theta}$ are constant.
\end{lemma}

\begin{proof}
A flip exchanges one ascending edge for another and retains the long edge $uw$.
Both short edges join $v$ to the component containing $u$ and $w$.
A corner is not contained in a larger natural edge, so it is never a short edge
of a flip.
\end{proof}

\subsection{Circular graphs}\label{subs:circularDefinitions}

A \defin{circular area sequence} is a tuple $\avec=(a_1,\dotsc,a_n)$ with
$n\geq1$, $0\leq a_i\leq n-1$ and $a_{i+1}\leq a_i+1$ for all $i$, with indices
read modulo $n$.
We extend $\avec$ periodically to all of $\setZ$ and put $b(i)=i-a_i$, so that
$b$ is weakly increasing and $b(i+n)=b(i)+n$.
The \defin{lifted natural edges} are the pairs $uv$ of integers with
$b(v)\leq u<v$.
They are invariant under translation by $n$, and their classes modulo $n$ are
the \defin{natural edges} of the \defin{circular unit-interval graph}.
A natural edge from the residue of $u$ to that of $v$ has length $v-u<n$; two
residues may be joined by natural edges in both directions, and we regard these
as different edges.
Corners and strict edges are defined on the lift as before, and give a
\defin{circular decorated unit-interval graph} $\Gamma=(\avec,S)$.
An area sequence with $a_1=0$ has no natural edge ending at $1$, and then
$a_v\leq v-1$ for all $v$; thus the ordinary graphs are precisely the circular
graphs in which no lifted natural edge crosses from one period to the next.
Circular unit-interval graphs and their chromatic quasisymmetric functions were
studied in~\cite{Ellzey2016,Ellzey2017,AlexanderssonPanova2018}.

Colorings $\kappa:\setZ/n\setZ\to\setNN$, ascents and $\LLT_\Gamma(\xvec;q)$ are
defined as before, with one ascent for every normal edge whose tail has a
smaller color than its head.
An orientation $\theta$ is again a set of ascending normal edges, and
$\Dplus{\theta}$ is the directed multigraph on $\setZ/n\setZ$ with the strict
and ascending edges.
Now $\Dplus{\theta}$ may contain directed cycles, for example around the whole
circle.
We call $\theta$ \defin{admissible} if $\Dplus{\theta}$ is acyclic, and we
write $\ORIENTS(\Gamma)$ for the set of admissible orientations.
A directed cycle admits no strictly increasing coloring, so
Lemma~\ref{lem:stateExpansion} holds verbatim, with the sum over admissible
orientations.

Let $\theta$ be admissible and let $\widetilde D_{+}(\theta)$ be the periodic
graph on $\setZ$ consisting of the lifted strict and ascending edges.
Its edges point upwards, so reachability in $\widetilde D_{+}(\theta)$ is a
partial order on $\setZ$, and we take covers in this order.
A path in $\widetilde D_{+}(\theta)$ cannot visit two lifts of the same residue,
since its projection would contain a directed cycle.
Hence every integer $z$ has a largest reachable integer
$\widetilde{\hrvop}(\theta,z)$, and
$\widetilde{\hrvop}(\theta,z+n)=\widetilde{\hrvop}(\theta,z)+n$.
The highest reachable vertex $\hrv{\theta}{u}$ is the residue of
$\widetilde{\hrvop}(\theta,u)$, and $\lambda(\theta)$ is the partition of the
fibre sizes of $\hrv{\theta}{\cdot}$ as before.
For ordinary graphs these definitions agree with the previous ones.

\begin{example}\label{ex:threeCycle}
Let $\avec=(1,1,1)$ with no strict edges.
The natural edges form the directed cycle $1\to2\to3\to1$.
The orientation with ascending edges $3\to1$ and $1\to2$ lifts to the path
$3\to4\to5$, so all three vertices have highest reachable vertex $2$.
The admissible orientations are the empty one, the three single edges and the
three paths of two edges, and the right-hand side of
Corollary~\ref{cor:circularExpansion} below is
$\elementaryE_{111}+3q\,\elementaryE_{21}+3q^2\elementaryE_3$.
\end{example}

Flips are defined on the lift.
Let $u<v<w<u+n$ be integers such that $uw$ is a lifted natural edge, and let
$\theta,\theta'$ be admissible orientations related as in
Definition~\ref{def:flip}, exchanging the residue classes of $uv$ and $vw$.
They differ by a \defin{flip} if the four cover conditions of
Definition~\ref{def:flip} hold in the reachability orders of
$\widetilde D_{+}(\theta)$ and $\widetilde D_{+}(\theta')$.
Lemma~\ref{lem:flipInvariants} holds for circular graphs as well, with the same
proof.

\subsection{Main results}\label{subs:mainResults}

For a set $X$ of orientations of $\Gamma$ we write
\[
 \defin{\Kv(X)}\coloneqq\sum_{\theta\in X}\Kp_{\Pplus{\theta}}(\xvec),
 \qquad
 \defin{\Ev(X)}\coloneqq\sum_{\theta\in X}\elementaryE_{\lambda(\theta)}(\xvec).
\]

\begin{theorem}\label{thm:main}
Let $\Gamma$ be an ordinary or circular decorated unit-interval graph.
Then $\Kv(C)=\Ev(C)$ for every flip class $C$ of $\Gamma$.
\end{theorem}

The number of ascending edges is constant on $C$ by
Lemma~\ref{lem:flipInvariants}; we denote it by $\asc(C)$.
Combining Theorem~\ref{thm:main} with Lemma~\ref{lem:stateExpansion}, every flip
class contributes
\[
 q^{\asc(C)}\Kv(C)=q^{\asc(C)}\sum_{\theta\in C}\elementaryE_{\lambda(\theta)}(\xvec)
\]
to $\LLT_\Gamma(\xvec;q+1)$, which is elementary-positive.
Summing over all classes we obtain the following.

\begin{corollary}\label{cor:circularExpansion}
For every circular decorated unit-interval graph $\Gamma$,
\[
 \LLT_\Gamma(\xvec;q+1)=\sum_{\theta\in\ORIENTS(\Gamma)}q^{\asc(\theta)}
 \elementaryE_{\lambda(\theta)}(\xvec).
\]
In particular, $\LLT_\Gamma(\xvec;q+1)$ is elementary-positive.
\end{corollary}

For ordinary graphs this is Theorem~\ref{thm:AS}.
Our proof does not use the recursion of~\cite{AlexanderssonSulzgruber2022}, and
thus gives a new proof of Theorem~\ref{thm:AS}.
